\documentclass[11pt,a4paper]{article}
\usepackage{amsmath,amssymb}
\newcommand{\Tr}{\operatorname{Tr}}
\begin{document}
\section*{Readout Jacobian}
Suppose we take a closed finite-dimensional system, with initial state $\rho$ and evolution $U_t=\exp(-itH/\hbar)$. Instead of comparing the density matrices directly, we take fixed Hermitian observables $A_a$ and study their expectation values. These are expected measurement readouts, not individual outcomes; we do not measure during the evolution.
\[x_a(t)=\Tr(A_aU_t\rho U_t^\dagger).\]
Now we perturb the initial state in directions specified by Hermitian generators $B_b$,
\[\rho_{\boldsymbol\epsilon}=V_{\boldsymbol\epsilon}\rho V_{\boldsymbol\epsilon}^\dagger,\qquad V_{\boldsymbol\epsilon}=\exp\!\left(-\frac{i}{\hbar}\sum_b\epsilon_bB_b\right).\]
We evolve each preparation with the same $U_t$. The Jacobian tells us how much each readout changes when we vary the initial preparation:
\[R_{ab}(t)=\left.\frac{\partial x_a(t,\boldsymbol\epsilon)}{\partial\epsilon_b}\right|_{\boldsymbol\epsilon=0}=\frac{i}{\hbar}\Tr(\rho[B_b,A_a(t)]),\qquad A_a(t)=U_t^\dagger A_aU_t.\]
Thus, to first order, we have
\[\delta\boldsymbol x(t)=R(t)\delta\boldsymbol\epsilon,\qquad\dot{\delta\boldsymbol x}(t)=\dot R(t)\delta\boldsymbol\epsilon.\]
If $R(t)$ is square and invertible, we can eliminate the initial parameters:
\[\dot{\delta\boldsymbol x}(t)=\dot R(t)R(t)^{-1}\delta\boldsymbol x(t).\]
However, this is only the response equation along our chosen preparation family. To obtain a closed evolution of $x$ for all states, we need the commutators to stay in the span of the readouts and the identity:
\[\frac{i}{\hbar}[H,A_a]=c_aI+\sum_bL_{ab}A_b\quad\Longrightarrow\quad\dot x=c+Lx,\qquad\dot{\delta x}=L\delta x.\]
The singular values of $R(t)$ measure response for specified input and output norms. They are not automatically Lyapunov growth factors: $R(0)$ need not be the identity, and closure does not itself imply instability.
\end{document}
